\subsection{A whole-residual complex child}
\label{sec:bulk-complex-child}

The complex phase interface admits a further batching step. On an edge
with residual dimension $a$, the existing binary change of basis puts
an orthonormal residual basis into the first $a$ slots. If each slot has
$f$ selected axes, those first $a$ slots form one contiguous interval
of $af$ axes. Instead of $a$ child calls on $f$ axes, use one child
on that interval.

The directions need not all have the same sign. A basis vector's weight
modulo four may make its directional kernel an inverse even on an
increasing edge. Let $B$ be the set of negative residual slots and
let $Z_B$ multiply an address by $-1$ precisely when the parity of its
selected bits in those slots is odd. The exact identity
$C^{-1}=-iZCZ$ gives
\[
 \bigotimes_{j=1}^{a}(C^{\varepsilon_j})^{\otimes f}
 =(-i)^{f|B|}Z_B C^{\otimes af}Z_B,
 \qquad \varepsilon_j\in\{1,-1\}.
\]
Thus one forward child, two address-sign passes and one fixed
fourth-root phase implement the mixed signs. The same binary basis
changes before and after the child are retained. The wrappers cost
$O(V)$: their sign is computed from the current address descriptor,
whose polynomial work per record is absorbed by the existing
superpolynomial record length. They introduce at most four elementary
coefficient operations per value and no further recursive call.

We batch only the two specified auxiliary edge classes of ranks
\[
 a_1=m-2h,\qquad a_2=m-h^2.
\]
All other edges continue to use their old one-direction children.
The first class is the increasing join between a completed stage-1
auxiliary bank, at dimension $h$, and its matched stage-3 bank, at
dimension $m-h$. The second is the stage-2 auxiliary exit from
dimension $h^2$ to the full terminal frame. These are complete
auxiliary banks, including the central roles when those are shared.
If a width bound requires padding, add dummy auxiliary roles with
identity scalar gates at the ordinary low and high frames before
matching the first- and third-stage banks. Their labels only increase,
their arbitrary scalar values are restored, and both counted auxiliary
classes are then present for every padded role. The changed
volume-weighted time expression is consequently
\[
 \frac1W\left[
    \left(s-n_1a_1-n_2a_2\right)m^{-\sigma}
    +n_1\left(\frac{a_1}{m}\right)^\sigma
    +n_2\left(\frac{a_2}{m}\right)^\sigma
 \right],
\]
where $n_1,n_2$ count the chosen edges. It has value $s/(Wm)<1$
at $\sigma=1$. Strict comparison below one supplies the new complex
recursion exponent. The separate mixed-width row construction handles
$f=\lfloor e/m\rfloor$, the at most $m-1$ individually processed
remaining axes, and a row reservation sufficient for the longer but
still $O(\log d)$ recursion.

\subsection{A depth guard for the unequal complex children}
\label{sec:bulk-complex-guard}

The dependency-path argument of Section~\ref{sec:pathwise-guard}
still bounds the sum of edge ranks on one scalar path by
\[
 q=m+6h.
\]
For $h=28$ its relevant constants are
\[
 m=21952,\quad q=22120,\quad a_1=21896,\quad a_2=21168.
\]
Both selected ranks exceed $q/2$. Therefore every arithmetic dependency
path encounters at most one selected bulk edge. At a call on $e$ axes,
put $f=\lfloor e/m\rfloor$. Ignoring the local additive operations,
the path has one of the following child-depth bounds:
\[
 qA(f),\qquad
 (q-a_1)A(f)+A(a_1f),\qquad
 (q-a_2)A(f)+A(a_2f).
\]
The weight budget accounts for both increasing and decreasing frame
edges and for paths that switch wires at a scalar gate. Inverse-child
and endpoint sign corrections add no child. Individually processing
the fewer than $m$ leftover axes adds at most $8m$ elementary
operations, absorbed into the local constant below.

Take $\rho=6/5$. The following three normalized path expressions have
a rational common upper bound strictly below one:
\begin{equation}\label{eq:bulk-path-moments}
 \max\left\{\frac q{m^\rho},\quad
 \frac{q-a_1}{m^\rho}+\left(\frac{a_1}m\right)^\rho,\quad
 \frac{q-a_2}{m^\rho}+\left(\frac{a_2}m\right)^\rho\right\}
 <\overline\theta:=\frac{999}{1000}.
\end{equation}
Here is an entirely rational verification. First
$160000^5<m^6$, so $m^\rho>160000$. For $0<t<1$, Taylor's formula
with its second derivative bounded on $[0,t]$ gives
\[
 (1-t)^{6/5}
 \le1-\frac65t+\frac3{25}\frac{t^2}{1-t}.
\]
Indeed its second derivative is
$(6/25)(1-t)^{-4/5}\le(6/25)/(1-t)$.
Putting $t=(m-a_i)/m$ gives upper bounds, in the displayed order,
\[
 \frac{553}{4000},\qquad
 \frac{47817969}{47897500},\qquad
 \frac{1213697}{1260000},
\]
all strictly below $999/1000$.

\paragraph{Stopping and exact arithmetic.}
Retain a global stopping threshold $d^\beta$, with $0<\beta<1$.
For sufficiently large $d$, this threshold exceeds $2m$. A leaf reached
from an internal call has width at least $d^\beta/(2m)$: every child
has width at least $f=\lfloor e/m\rfloor\ge e/(2m)$. Put
\[
 E=64(W+m+1)^3,
 \qquad C_{\rm dep}=1000(E+16m+1).
\]
The exact finite inequality
$36W^3+4s+4W+8m+4<E$ covers all local scalar operations, the new
leftover-axis work and the retained endpoint and inverse wrappers.
For every descendant width $e$ reached before or just after stopping,
strong induction gives
\begin{equation}\label{eq:bulk-depth-power}
 A(e)\le C_{\rm dep}d^{-\beta(\rho-1)}e^\rho.
\end{equation}
At a leaf, use
\[
 8e\le8(2m)^{\rho-1}d^{-\beta(\rho-1)}e^\rho
       \le16m\,d^{-\beta(\rho-1)}e^\rho.
\]
At an internal call, child widths are at most their stated fractions
of $e$, so \eqref{eq:bulk-path-moments} bounds their total by
$\overline\theta C_{\rm dep}d^{-\beta(\rho-1)}e^\rho$.
Since $e\ge d^\beta$, the remaining gap is at least
$(1-\overline\theta)C_{\rm dep}d^\beta\ge E$. This proves the induction.
A root that is already a leaf separately satisfies $A(e)\le8d^\beta$,
which is smaller than the final bound below.

For rational $\zeta>0$, set
\[
 C_1=\rho-(\rho-1)\beta+\zeta,
 \qquad
 C_0=\left\lceil128m(1+1/\zeta)C_{\rm dep}\right\rceil,
 \qquad \Delta=\lceil C_0d^{C_1}\rceil.
\]
The root bound from \eqref{eq:bulk-depth-power} is
$C_{\rm dep}d^{\rho-(\rho-1)\beta}$. This covers the direct arbitrary-width
root construction; it also covers every retained base-$m$ piece.
Even retaining the older bound of
$m(1+1/\zeta)d^\zeta$ such pieces, their concatenation, all reserved
axis kernels, the outer phases and the $b^D$ conversion have depth
at most $C_0d^{C_1}$. The original elementary-depth induction
therefore supplies exact denominator and magnitude guards of
$p+\Delta$ fractional bits and $\Delta+3$ integer/sign bits.
The width remains $O(p)$ when $\epsilon C_1<1$.

In particular, $\beta=1/1000$ and $\zeta=1/10000$ give
\[
 C_1=\frac{11999}{10000},\qquad
 1-\frac{7999999}{16000000}C_1
   =\frac{64008011999}{160000000000}>0.
\]
Thus bulk complex children preserve the near-half dimension exponent
while allowing the improved complex time exponent to be used with
almost no loss at the stopping leaves.
